\chapter{Tutorials}
\label{ch:tutorials_pyclical}

\section{Notation and Basic Operations}
PyClical uses a specific notation for Clifford algebra elements.
The \texttt{clifford} class represents multivectors, and the \texttt{index\_set} class represents the basis elements.

\subsection{Creating Multivectors}
You can create multivectors using the \texttt{e()} helper function (short for \texttt{clifford(index\_set(...))}).
\begin{verbatim}
>>> x = 2*e(1) + 3*e(2)
>>> print(x)
2{1}+3{2}
\end{verbatim}
Generators with negative indices square to $-1$, and those with positive indices square to $+1$.

\section{Rotations in $\Real^3$}
Rotations can be performed using rotors (exponentials of bivectors).
\begin{verbatim}
>>> x = 2*e(1)+3*e(2)+5*e(3)
>>> a = (3*e(1)+2*e(2)+e(3))/10
>>> s = exp(e({1,2,3})*a/2)
>>> y = s*x/s
>>> print(y*y)
38
\end{verbatim}
The norm is preserved.

\section{Spacetime Physics}
PyClical can model Minkowski spacetime $\Real^{3,1}$.
\begin{verbatim}
>>> x = 2*e(-1)+2*e(1)+3*e(2)+5*e(3)
\end{verbatim}
Lorentz transformations are implemented as rotors in this algebra.

\section{Advanced Topics}
PyClical is versatile enough to handle:
\begin{itemize}
    \item Octonions (via Cayley algebra construction).
    \item Conformal Geometric Algebra (CGA).
    \item Functions of multivectors (sqrt, log, exp).
\end{itemize}
See the \texttt{pyclical\_demo.py} and \texttt{sqrt\_log\_demo.py} files for advanced examples.
